\section*{Capítulo \ref{ch:b1:observing-the-sky} --- Telescópios, detectores e fotometria}

\begin{solution}{exo:b1:observing-the-sky:1}
Uma diferença de 5 magnitudes corresponde a um fator de 100: a estrela de magnitude 1.0 envia 100 vezes mais \omterm{def:b1:observing-the-sky:flux}{fluxo}. Rigel e Polaris: $10^{0.4\times(2.02 - 0.13)} =
10^{0.756} = 5.70$; Rigel, de magnitude menor, é mais brilhante.
\end{solution}

\begin{solution}{exo:b1:observing-the-sky:2}
$d = 1/0.19495 = \qty{5.13}{\parsec}$, isto é, $5.1295\times206265 =
\qty{1.058e6}{\au}$ e $5.1295\times3.2616 = \qty{16.7}{\lightyear}$.
\end{solution}

\begin{solution}{exo:b1:observing-the-sky:3}
$1.22\times\num{550e-9}/2.4 = \qty{2.80e-7}{rad} = \ang{;;0.058}$, cerca de
sete vezes mais fino que um \omterm{def:b1:observing-the-sky:seeing}{seeing} de \ang{;;0.40}: no espaço o telescópio
atinge seu limite de difração. Para \qty{6.5}{m} em \qty{2.0}{\micro\metre}:
$1.22\times\num{2.0e-6}/6.5 = \qty{3.75e-7}{rad} = \ang{;;0.077}$.
\end{solution}

\begin{solution}{exo:b1:observing-the-sky:4}
$B - V$ cresce de Rigel ($-0.03$) para Altair (0.22), Polaris (0.60) e
Betelgeuse (1.85): a mais quente é Rigel, depois Altair, Polaris e
Betelgeuse. A distância reduz as duas bandas pelo mesmo fator, de modo que $B - V$
não é alterado por ela (desde que nenhuma poeira avermelhe a luz no caminho).
\end{solution}

\begin{solution}{exo:b1:observing-the-sky:5}
$X_1 = \sec\ang{30} = 1.1547$, $X_2 = \sec\ang{60} = 2.0000$.
$k = (6.590 - 6.412)/(2.0000 - 1.1547) = 0.178/0.8453 = \num{0.211}$ por
\omterm{def:b1:observing-the-sky:extinction}{massa de ar}, e $m_0 = 6.412 - 0.2106\times1.1547 = 6.169$.
\end{solution}

\begin{solution}{exo:b1:observing-the-sky:6}
$\mathrm{SNR} = \sqrt{400} = 20$. A SNR cresce como $\sqrt t$: para 100,
$t = \qty{10}{s}\times(100/20)^2 = \qty{250}{s}$. Um erro de \num{0.01}
magnitude exige $\mathrm{SNR} = 1.086/0.01 = 108.6$, logo
$t = 10\times(108.6/20)^2 = \qty{295}{s}$.
\end{solution}

\begin{solution}{exo:b1:observing-the-sky:7}
A distância é $d = 1/0.00378 = \qty{264.6}{\parsec}$, e a \omterm{def:b1:observing-the-sky:absolute}{magnitude
absoluta} $M = 0.13 - 5\log_{10}26.46 = 0.13 - 7.11 = -6.98$. $L/L_\odot = 10^{0.4\times(4.83 + 6.98)} =
10^{4.724} = \num{5.3e4}$ na banda V.
\end{solution}

\begin{solution}{exo:b1:observing-the-sky:8}
$\pi\times1.2^2 = \qty{4.52}{m^2}$, e $25/4.52 = 5.5$: o telescópio Webb
coleta o mesmo número de fótons em uma exposição 5.5 vezes mais curta. Um
disco cheio de \qty{6.5}{m} de diâmetro teria $\pi\times3.25^2 = \qty{33.2}{m^2}$;
o espelho real tem $25/33.2 = \qty{75}{\percent}$ dele, porque 18 hexágonos
deixam descoberta a borda desse disco, e as frestas entre os segmentos e a
abertura central não coletam nada.
\end{solution}

\begin{solution}{exo:b1:observing-the-sky:9}
(a) $\qty{10}{\parsec} = \qty{3.0857e17}{m}$, e $F = \num{3.828e26}/(4\pi
\times(\num{3.0857e17})^2) = \num{3.828e26}/\num{1.1965e36} =
\qty{3.20e-10}{W.m^{-2}}$, $1361/\num{3.20e-10} = \num{4.3e12}$ vezes menos do que na
Terra.
(b) $5\log_{10}(d/\qty{10}{\parsec}) = -26.74 - 4.83 = -31.57$, logo
$d = \qty{10}{\parsec}\times10^{-6.314} = \qty{4.853e-6}{\parsec}$, isto é,
$\num{4.853e-6}\times206265 = \qty{1.001}{\au}$.
\end{solution}

\begin{solution}{exo:b1:observing-the-sky:10}
Acima da atmosfera o \omterm{def:b1:observing-the-sky:flux}{fluxo} é $F_0$; abaixo, $F = F_0e^{-\tau_0X}$. Então
$m - m_0 = -2.5\log_{10}(F/F_0) = 2.5\log_{10}e\;\tau_0X$, que é $kX$ com
$k = 2.5\log_{10}e\;\tau_0 = 1.0857\,\tau_0$. Para $k = 0.160$:
$\tau_0 = 0.160/1.0857 = 0.147$. No \omterm{def:b1:spherical-astronomy:celestial-sphere}{zênite} chega $e^{-0.147} = 0.863$ da
luz, \qty{86}{\percent}; em $X = 2.31$, $e^{-0.147\times2.31} =
e^{-0.340} = 0.712$, \qty{71}{\percent}. Em $z = \ang{85}$, $\sec z =
11.5$; mas o modelo plano-paralelo torna o percurso infinito no horizonte,
ao passo que o percurso por uma atmosfera curva permanece finito, e a refração
também desvia o raio: perto do horizonte a fórmula superestima a
\omterm{def:b1:observing-the-sky:extinction}{massa de ar}, e seus erros crescem depressa.
\end{solution}

\begin{solution}{exo:b1:observing-the-sky:11}
(a) $\mathrm{SNR} = Rt/\sqrt{(R + B)t}$, logo $t = \mathrm{SNR}^2(R + B)/R^2
= 100\times2050/2500 = \qty{82}{s}$. (b) Reduzir o \omterm{def:b1:observing-the-sky:seeing}{seeing} à metade divide a
área da imagem, e portanto $B$, por 4: $B = 500$ e $t = 100\times550/2500 =
\qty{22}{s}$.
(c) Para $B \gg R$, $t \approx \mathrm{SNR}^2B/R^2$. A taxa $R$ da fonte cresce como $D^2$;
a taxa $B$ do céu cresce como a área coletora, $D^2$, vezes a área do céu na imagem,
$\theta^2$. Logo $t \propto
\theta^2D^2/D^4 = \theta^2/D^2$: reduzir o \omterm{def:b1:observing-the-sky:seeing}{seeing} à metade vale tanto quanto
dobrar o diâmetro.
\end{solution}

\begin{solution}{exo:b1:observing-the-sky:12}
(a) $N = 1200\times100 = \num{120000}$ linhas; na primeira ordem o poder de
resolução é \num{1.2e5}, e $\Delta\lambda = 550/(\num{1.2e5}) =
\qty{4.6e-3}{nm}$. (b) $v = c\,\Delta\lambda/\lambda = c/(\num{1.2e5}) =
\qty{2.5}{km.s^{-1}}$. (c) $\Delta\lambda = 550\times29.78/\num{299792} =
\qty{0.055}{nm}$, cerca de 12 elementos de resolução. O observador deve calcular,
para a data e a hora da observação, a componente da velocidade orbital da Terra
na direção da estrela, e subtraí-la, para obter a velocidade da estrela em
relação ao Sol.
\end{solution}

\begin{solution}{pb:b1:observing-the-sky:1}
\textbf{1.} $d = 1/0.37921 = \qty{2.637}{\parsec}$.

\textbf{2.} $2.637\times206265 = \qty{5.44e5}{\au}$;
$2.637\times\num{3.0857e16} = \qty{8.14e16}{m}$.

\textbf{3.} $2.637\times3.2616 = \qty{8.60}{\lightyear}$: a luz que vemos saiu de
Sirius há 8.6 anos.

\textbf{4.} A direção oscila de $2p = \ang{;;0.758}$, cerca do dobro do
\omterm{def:b1:observing-the-sky:seeing}{seeing}. A oscilação é medida a partir da posição do centro da imagem, e não da sua borda: em média sobre muitos fótons,
o centro é encontrado com uma pequena fração da largura da imagem, tanto mais
precisamente quanto mais fótons forem contados.

\textbf{5.} $\Delta d/d = \Delta p/p = 0.001/0.37921 = \qty{0.26}{\percent}$,
isto é, $8.60\times0.0026 = \qty{0.023}{\lightyear}$.

\textbf{6.} $10^{0.4\times(0.03 + 1.46)} = 10^{0.596} = 3.94$.

\textbf{7.} $10^{0.4\times(26.74 - 1.46)} = 10^{10.112} = \num{1.29e10}$.

\textbf{8.} $10^{0.4\times(8.44 + 1.46)} = 10^{3.96} = \num{9120}$.

\textbf{9.} $z = \ang{90} - \ang{82.09} = \ang{7.91}$, $X = \sec\ang{7.91} =
1.0096$; extinção $0.15\times1.0096 = 0.151$ magnitude; medida,
$m = -1.46 + 0.151 = -1.31$.

\textbf{10.} $z = \ang{60}$, $X = 2.000$; extinção de 0.30 magnitude,
$m = -1.16$. A fração transmitida é $10^{-0.4\times0.30} = 0.759$:
\qty{24}{\percent} da luz se perde.

\textbf{11.} Ao nascer $z = \ang{90}$ e $\sec z$ é infinito, enquanto o
percurso real pela atmosfera curva é longo mas finito; a refração também
desvia o raio. Perto do horizonte a extinção é grande e muda depressa, e a
correção plano-paralela está errada.

\textbf{12.} $\num{3.0e4}\times10^{0.4\times(10.0 - 8.44)} = \num{3.0e4}
\times10^{0.624} = \num{3.0e4}\times4.21 = \num{1.26e5}$ fotoelétrons por
segundo.

\textbf{13.} Uma SNR de 1000 exige $N = 1000^2 = \num{1.0e6}$
fotoelétrons: $t = \num{1.0e6}/\num{1.26e5} = \qty{7.9}{s}$. O erro na
magnitude é $1.086/1000 = \num{0.0011}$.

\textbf{14.} $\num{3.0e4}\times10^{0.4\times(10.0 + 1.46)} = \num{3.0e4}
\times10^{4.584} = \num{1.15e9}$ fotoelétrons por segundo.

\textbf{15.} $0.20\times\num{1.15e9} = \num{2.30e8}$ por segundo no
pixel mais brilhante, que se enche em $\num{1.0e5}/\num{2.30e8} =
\qty{4.3e-4}{s}$, menos de meio milissegundo.

\textbf{16.} As asas de Sirius dão $\num{1.15e9}\times10^{-4} =
\num{1.15e5}$ por segundo; Sirius~B dá $0.20\times\num{1.26e5} =
\num{2.52e4}$ por segundo. O ofuscamento de Sirius é 4.6 vezes mais brilhante que
o pico de Sirius~B.

\textbf{17.} O pixel recebe $\num{1.15e5} + \num{2.52e4} = \num{1.40e5}$
por segundo e se enche em $\num{1.0e5}/\num{1.40e5} = \qty{0.71}{s}$. Em
\qty{0.5}{s}: sinal $\num{2.52e4}\times0.5 = \num{1.26e4}$, total
$\num{7.01e4}$, ruído $\sqrt{\num{7.01e4}} = 265$, logo $\mathrm{SNR} =
\num{1.26e4}/265 = 48$. Sirius~B é detectável do ponto de vista do ruído de
fótons; o que a torna difícil é que o ofuscamento precisa ser conhecido e
subtraído com precisão muito melhor que um quinto do seu valor, enquanto o núcleo
de Sirius satura em menos de um milissegundo.

\textbf{18.} $m - M = 5\log_{10}(2.637/10) = 5\times(-0.579) = -2.89$.

\textbf{19.} $M = -1.46 + 2.89 = 1.43$.

\textbf{20.} $d = 1/0.13023 = \qty{7.679}{\parsec}$;
$m - M = 5\log_{10}0.7679 = -0.57$, logo $M = 0.03 + 0.57 = 0.60$.

\textbf{21.} Vega, de \omterm{def:b1:observing-the-sky:absolute}{magnitude absoluta} menor, por um fator de
$10^{0.4\times(1.43 - 0.60)} = 10^{0.332} = 2.15$. Sirius parece mais brilhante
apenas porque está mais perto.

\textbf{22.} $m = M_\odot + (m - M) = 4.83 - 2.89 = 1.94$: o Sol seria
uma estrela bem brilhante a olho nu.

\textbf{23.} $5\log_{10}(d/\qty{10}{\parsec}) = 6.0 - 1.43 = 4.57$, logo
$d = 10\times10^{0.914} = \qty{82.0}{\parsec}$, cerca de \qty{267}{\lightyear}.

\textbf{24.} $L/L_\odot = 10^{0.4\times(4.83 - 1.43)} = 10^{1.36} = 22.9$:
na banda V, Sirius ofusca \textbf{cerca de 23 sóis}.
\end{solution}
